\section{과업 수준 준거 연관성 결과}
\label{sec:validity-results}

\subsection{Development freeze}

첫 development specification은 sealed reveal 전에 실패했다.
Leave-one-world-out improvement는 terminal Brier score에서 \(0.000421\),
integrated score에서 \(-0.000221\)이었다. Effect floor를 낮추지 않았다.
Revised specification은 exogenous task utility를 mode equality에서 structural
predicate로 바꾸고, task-varying hazard와 scalar
selected-versus-alternative contrast를 사용했으며, world당 unit을 16에서
32로 늘리고 stronger layer-aware sensitivity를 유지했다.

Revised development world 24개에서 primary event rate는 \(0.60384\)였다.
Leave-one-world-out terminal/integrated improvement는 각각 \(0.003124\),
\(0.003927\)이었다. Analytic floor와 20,000-draw conjunctive Holm power
simulation으로 sealed world 69개를 선택했다. Simulated conjunctive power는
\(0.9208\), Monte Carlo 95\% lower bound는 \(0.91698\)이었다.

\subsection{Sealed primary analysis}

Sealed fit \(1{,}242\)개가 모두 완료됐다. Primary population은
\(26{,}496=69\times32\times3\times4\)개 unit record(world, unit, nested
optimizer seed, primary arm)를 포함했고 event rate는 \(0.61632\)였다. 두 exact
control의 event rate는 0을 유지했다. Table~\ref{tab:validity-primary}는
co-primary analysis를 보인다.

\begin{table}[ht]
\centering
\caption{69개 world의 frozen co-primary criterion-validity 결과. Bootstrap과
hierarchical column은 0에 대한 simultaneous lower bound다.}
\label{tab:validity-primary}
\begin{tabular}{lrrrrr}
\toprule
Effect & Mean & World SD & Holm \(p\) & Bootstrap LB & Hierarchical LB \\
\midrule
Terminal Brier
& 0.003519 & 0.009015 & 0.001230 & 0.001428 & 0.001392 \\
Integrated first failure
& 0.003034 & 0.007471 & 0.001230 & 0.001296 & 0.001271 \\
\bottomrule
\end{tabular}
\end{table}

두 mean은 \(\SESOI=0.0025\)를 넘는다. Mean terminal Brier score는
\(0.183661\)에서 \(0.180143\)으로, integrated Brier score는 \(0.148317\)에서
\(0.145283\)으로 감소했다. Outcome noncircularity, panel completeness, event
balance, no-refitting, exact-control structural zero, Holm,
cluster-bootstrap, hierarchical requirement가 모두 통과했다.

Pooled estimand는 이질적인 네 arm을 평균한다. 따라서 routing interpretation에
앞서 사전지정 secondary arm decomposition을
Table~\ref{tab:validity-arms}에 보고한다.

\begin{table}[ht]
\centering
\caption{Frozen scalar model 대비 arm-level Brier improvement. 별도로
multiplicity gate를 적용한 estimand가 아니라 secondary descriptive effect다.}
\label{tab:validity-arms}
\begin{tabular}{lrr}
\toprule
Arm & Terminal & Integrated \\
\midrule
Layer-routed dense & 0.005402 & 0.003222 \\
Strict factorized & 0.003114 & 0.003077 \\
Random-routed matched sparsity & 0.002445 & 0.002758 \\
Permuted or wrong-routed & 0.003114 & 0.003077 \\
\bottomrule
\end{tabular}
\end{table}

Wrong-routed arm은 두 outcome 모두에서 \(\SESOI\)를 넘고 random-routed arm은
integrated outcome에서 이를 넘는다. 따라서 pooled criterion analysis는
structural augmentation과 same-task outcome 사이의 association을 식별하지만
routing correctness의 효과를 식별하지 않는다. Layer-routed dense의 descriptive
effect가 strict factorization보다 크다는 사실도 factorization을 작동 성분으로
분리하지 못한다.

\subsection{Sensitivity와 temporal boundary}

Generic baseline에 direct raw layer mismatch를 추가했을 때 structural
augmentation의 terminal/integrated Brier improvement는 \(0.002108\),
\(0.001467\)이었다. 양수이지만 \(\SESOI\) 미만이다. 따라서 primary result는
scalar generic baseline을 넘는 increment를 확립하지만 prespecified
layer-aware baseline을 넘는 practically nontrivial advantage는 확립하지
않는다.

더 강한 temporal interpretation 두 개를 development data에서만 검사했으며
모두 실패했다. Independent probe로 later task failure를 예측했을 때
\(-0.000539\), \(-0.000212\)였고, 처음 네 task를 calibration으로 사용해 뒤
네 task를 예측했을 때 \(-0.002021\), \(-0.001069\)였다. 따라서 sealed
positive result는 oracle candidate endpoint를 사용하는 same-task criterion
association이다. Future-task prognosis 또는 oracle-free deployment score의
evidence가 아니다.
